% !TeX root = ../../Book3.tex
\section{\texorpdfstring{\(A^2\)}{A²} 的完备性与正交基}
\label{b3:sec:2804}
Hardy 范数限制每一条同心圆周上的平方平均，因而在半径趋于一时仍保留边界控制。若只要求整个圆盘上的平方积分有限，靠近边界的增长就可以更大。这个改动会把 Taylor 系数的普通平方和变成带权平方和，也会改变点值泛函的最佳常数。本节先从面积积分直接建立空间的完备性，再用圆盘的旋转对称性找出它的正交基。

\subsection{面积范数与局部控制}
\label{b3:subsec:280401}

\begin{definition}[圆盘 Bergman 空间]\label{b3:def:28-bergman-space}
设 \(\mathbb D=\{\,z\in\mathbb C\mid |z|<1\,\}\)，取归一化面积测度 \(\mathrm{d}A(z)=\mathrm{d}x\,\mathrm{d}y/\pi\)，其中 \(z=x+\mathrm{i}y\)。圆盘上的\term[bergman-kongjian]{Bergman 空间}[Bergman space]定义为
\[
 A^2(\mathbb D)\coloneqq
 \left\{\,f\in\mathcal H(\mathbb D)
 \;\middle|\;\int_{\mathbb D}|f(z)|^2\,\mathrm{d}A(z)<\infty\,\right\}.
\]
在这个空间上定义范数
\[
 \|f\|_{A^2}\coloneqq
 \left(\int_{\mathbb D}|f(z)|^2\,\mathrm{d}A(z)\right)^{1/2}.
\]
对 \(f,g\in A^2(\mathbb D)\)，规定
\[
 \langle f,g\rangle_{A^2}
 \coloneqq\int_{\mathbb D}f(z)\cdot\overline{g(z)}\,\mathrm{d}A(z).
\]
这个复内积对第一变量线性，对第二变量共轭线性。
\end{definition}
\symidx{\(A^2(\mathbb D)\)：圆盘 Bergman 空间}

归一化使 \(A(\mathbb D)=1\)，所以常数函数 \(1\) 的范数等于一。若改用未经归一化的面积 \(\mathrm{d}x\,\mathrm{d}y\)，范数将乘以 \(\sqrt\pi\)，后面的核则须除以 \(\pi\)。同一计算中的这两种约定不能混用。这里的面积范数也不同于 \(H^2\) 的边界范数，尽管它们都带上标二。

\ProseParagraphBreak

一般 \(L^2\) 函数只按几乎处处相等识别，改变单点的值不改变范数。全纯函数之间却没有这个含混：若全纯函数 \(f-g\) 在一个内点不为零，由连续性，它在某个小圆盘上的模便有正下界，因而不可能几乎处处为零。因此每个具有全纯代表的 \(L^2(\mathbb D,\mathrm{d}A)\) 等价类，至多有一个全纯代表；我们以后就使用这个代表。下面的估计进一步说明，其点值确实由面积范数连续控制。

\begin{proposition}[面积平均对点值的控制]\label{b3:prop:28-area-local-control}
设 \(f\) 在单位圆盘 \(\mathbb D\) 全纯，\(w\in\mathbb D\)，且 \(0<\rho<1-|w|\)。对于归一化面积测度 \(\mathrm{d}A=\mathrm{d}x\,\mathrm{d}y/\pi\)，有
\begin{equation}\label{b3:eq:28-area-mean-control}
 |f(w)|^2\le\frac1{\rho^2}
 \int_{D(w,\rho)}|f(z)|^2\,\mathrm{d}A(z).
\end{equation}
特别地，若 \(f\in A^2(\mathbb D)\)，则对每个 \(0\le r<1\)，
\[
 \sup_{|w|\le r}|f(w)|\le\frac{\|f\|_{A^2}}{1-r}.
\]
\end{proposition}
\begin{proof}
当 \(0<t<\rho\) 时，圆周 Cauchy 公式给出
\[
 f(w)=\frac1{2\pi}\int_0^{2\pi}
 f(w+te^{\mathrm{i}\theta})\,\mathrm{d}\theta.
\]
在总质量为一的圆周测度上使用 Cauchy--Schwarz 不等式，得到
\[
 |f(w)|^2\le\frac1{2\pi}\int_0^{2\pi}
 |f(w+te^{\mathrm{i}\theta})|^2\,\mathrm{d}\theta.
\]
两边乘以 \(2t\)，从 \(0\) 积到 \(\rho\)。左边为 \(\rho^2|f(w)|^2\)，右边由极坐标换元恰好等于归一化面积积分。这证明了\eqref{b3:eq:28-area-mean-control}。

若 \(|w|\le r\)，可以对全部这些 \(w\) 使用同一个 \(0<\rho<1-r\)，并把小圆盘积分放大为整个 \(\mathbb D\) 上的积分。于是 \(\sup_{|w|\le r}|f(w)|\le\|f\|_{A^2}/\rho\)。令 \(\rho\uparrow1-r\) 即得结论；无需假设 \(f\) 连续延拓到单位圆周。
\end{proof}

这个论证是\subsecref{b3:subsec:040302}所用局部平方平均估计的圆盘版本，这里已经按 \(\mathrm{d}A=\mathrm{d}x\,\mathrm{d}y/\pi\) 重写常数。它说明面积积分虽不单独限制一般可测函数的点值，却能限制全纯函数在每个固定内部紧集上的大小。全纯性通过圆周平均，把一个内点与周围一片正面积区域联系起来。

\ProseParagraphBreak

靠近边界时，上界允许增长。这个允许量有实际内容：\(\sqrt{n+1}\,z^n\) 的面积范数稍后将算出恒为一，但它在接近单位圆周处可以很大。在固定的 \(|z|\le r<1\) 上，\(\sqrt{n+1}\,r^n\to0\)，又与局部控制完全相容。下一节将把当前的粗常数改进到最佳值。

\subsection{Bergman 空间的完备性}
\label{b3:subsec:280402}

\(L^2\) 的完备性首先给出可测函数类的极限；要建立 \(A^2\) 的完备性，还须把这一极限识别为全纯函数。局部估计正好连接这两个空间。它将面积范数中的 Cauchy 条件变为每个闭子圆盘上的一致 Cauchy 条件；得到全纯极限以后，再回到面积积分证明范数收敛。

\begin{theorem}[圆盘 Bergman 空间的完备性]\label{b3:thm:28-bergman-complete}
在单位圆盘 \(\mathbb D\) 上取 \(\mathrm{d}A=\mathrm{d}x\,\mathrm{d}y/\pi\)。空间 \(A^2(\mathbb D)\) 关于内积 \(\langle f,g\rangle_{A^2}=\int_{\mathbb D}f\overline g\,\mathrm{d}A\) 是复 Hilbert 空间。以唯一的全纯代表将它视为 \(L^2(\mathbb D,\mathrm{d}A)\) 的子空间时，这个子空间是闭的。
\end{theorem}
\begin{proof}
线性组合仍全纯且平方可积，内积的各项性质来自积分；正定性由全纯代表的唯一性保证。只需证明完备性与闭性。

设 \((f_j)\) 是 \(A^2\) 中的 Cauchy 列。对固定 \(0<r<1\)，由\cref{b3:prop:28-area-local-control}作用于 \(f_j-f_k\)，有
\[
 \sup_{|z|\le r}|f_j(z)-f_k(z)|
 \le\frac1{1-r}\|f_j-f_k\|_{A^2}\longrightarrow0
 \quad(j,k\to\infty).
\]
所以每个闭子圆盘上都有一致极限。在两个闭子圆盘的重叠处，这些极限等于同一数列的逐点极限，因而彼此相容，给出 \(\mathbb D\) 上的函数 \(f\)。任何内部紧集都包含于某个闭子圆盘，故 \(f_j\to f\) 局部一致。各 \(f_j\) 全纯，\cref{b3:prop:04-holomorphic-limit}于是保证 \(f\) 全纯。

还须确认 \(f\) 具有有限的全盘面积范数，而不仅在每个紧集上良好。Cauchy 列有界，设 \(\|f_j\|_{A^2}\le M\)。对非负函数 \(|f_j|^2\) 使用 Fatou 引理，得到
\[
 \int_{\mathbb D}|f|^2\,\mathrm{d}A
 \le\liminf_{j\to\infty}\int_{\mathbb D}|f_j|^2\,\mathrm{d}A
 \le M^2,
\]
故 \(f\in A^2\)。给定 \(\varepsilon>0\)，选择 \(N\) 使 \(j,k\ge N\) 时 \(\|f_j-f_k\|_{A^2}<\varepsilon\)。固定 \(j\ge N\)，再让 \(k\to\infty\)，对 \(|f_j-f_k|^2\) 使用 Fatou 引理，得
\[
 \|f_j-f\|_{A^2}^2
 \le\liminf_{k\to\infty}\|f_j-f_k\|_{A^2}^2
 \le\varepsilon^2.
\]
因此整列在面积范数中趋于 \(f\)，完备性得证。

最后，若 \(f_j\in A^2\) 在 \(L^2(\mathbb D,\mathrm{d}A)\) 中趋于某个等价类 \(u\)，则它也是上述 Cauchy 列，故存在全纯 \(f\in A^2\) 使 \(\|f_j-f\|_{L^2}\to0\)。三角不等式给出 \(\|u-f\|_{L^2}=0\)，所以 \(u\) 正是这个全纯代表的等价类。这证明了子空间的闭性。
\end{proof}

这里两次使用 Fatou 引理的目的不同：第一次证明候选极限仍平方可积，第二次把 Cauchy 误差传给最终极限。我们没有从局部一致收敛直接断言全盘积分收敛，因为边界附近仍可能有质量集中。例如 \(\sqrt{n+1}\,z^n\) 局部一致趋于零，面积范数却不趋于零；它不是面积范数中的 Cauchy 列。

\ProseParagraphBreak

范数收敛还控制固定阶导数。若 \(f_j\to f\) 于 \(A^2\)，取 \(0<r<s<1\)，先在 \(|z|\le s\) 上使用局部估计，再在每个 \(|w|\le r\) 周围使用半径 \(s-r\) 的 Cauchy 导数估计，便有
\[
 \sup_{|w|\le r}|f_j^{(m)}(w)-f^{(m)}(w)|
 \le\frac{m!}{(s-r)^m(1-s)}\|f_j-f\|_{A^2}
 \quad(m=0,1,2,\ldots).
\]
这使面积范数中的极限与前文熟悉的全纯函数局部极限相容，但没有把任何结论延伸到移动的边界点上。

\subsection{单项式的正交性}
\label{b3:subsec:280403}

圆盘与面积测度都在旋转下不变。因此不同角频率应当互相正交；单项式 \(z^n=r^ne^{\mathrm{i}n\theta}\) 正好把角频率与径向幂次分开。直接计算还能看出，为什么面积空间的归一化基与 Hardy 空间不同。

\begin{proposition}[单项式的面积内积]\label{b3:prop:28-monomial-inner-product}
在 \(A^2(\mathbb D)\) 中取 \(\mathrm{d}A=\mathrm{d}x\,\mathrm{d}y/\pi\)。对任意非负整数 \(n,m\)，有
\[
 \langle z^n,z^m\rangle_{A^2}
 =\begin{cases}
 0,&n\ne m,\\
 \dfrac1{n+1},&n=m.
 \end{cases}
\]
因此函数 \(e_n(z)=\sqrt{n+1}\,z^n\)，\(n=0,1,2,\ldots\)，构成 \(A^2(\mathbb D)\) 的标准正交系。
\end{proposition}
\begin{proof}
记 \(\mathrm{d}m(\theta)=\mathrm{d}\theta/(2\pi)\)。极坐标面积元为 \(\mathrm{d}A=2r\,\mathrm{d}r\,\mathrm{d}m\)，所以
\[
 \langle z^n,z^m\rangle_{A^2}
 =2\int_0^1r^{n+m+1}\,\mathrm{d}r
 \int_0^{2\pi}e^{\mathrm{i}(n-m)\theta}\,\mathrm{d}m(\theta).
\]
当 \(n\ne m\) 时，角积分为零；当 \(n=m\) 时，角积分为一，而径向积分为 \(2/(2n+2)=1/(n+1)\)。乘上 \(\sqrt{n+1}\) 正好把每个单项式的长度变成一。
\end{proof}

对于多项式 \(p(z)=\sum_{n=0}^Na_nz^n\)，有限和的正交性立即给出
\[
 \|p\|_{A^2}^2=\sum_{n=0}^N\frac{|a_n|^2}{n+1}.
\]
例如 \(\|1+2z\|_{A^2}^2=1+4/2=3\)，而它的 Hardy 范数平方为 \(1+4=5\)。面积积分对高次项给出的权重较小，是因为 \(r^n\) 只有靠近边界时才不小；把全部半径再积一次，会减轻这一圈的贡献。

\ProseParagraphBreak

这些单项式相互正交，还需要知道它们能否逼近任意 \(A^2\) 函数。Taylor 展开与非负项积分将说明，多项式逼近可以从内部圆盘上的一致收敛提升为整个圆盘上的面积范数收敛。

\subsection{系数表示与正交基的完备性}
\label{b3:subsec:280404}

\begin{theorem}[Bergman 范数与 Taylor 系数]\label{b3:thm:28-bergman-coefficients}
设 \(f(z)=\sum_{n=0}^{\infty}a_nz^n\) 在单位圆盘 \(\mathbb D\) 全纯，并取 \(\mathrm{d}A=\mathrm{d}x\,\mathrm{d}y/\pi\)。则
\begin{equation}\label{b3:eq:28-bergman-coefficient-norm}
 \int_{\mathbb D}|f(z)|^2\,\mathrm{d}A(z)
 =\sum_{n=0}^{\infty}\frac{|a_n|^2}{n+1},
\end{equation}
允许两边同为 \(+\infty\)。特别地，\(f\in A^2(\mathbb D)\) 当且仅当右端有限；在此情形，Taylor 部分和 \(p_N(z)=\sum_{n=0}^Na_nz^n\) 在 \(A^2\) 范数中趋于 \(f\)，且
\[
 \|f-p_N\|_{A^2}^2=\sum_{n>N}\frac{|a_n|^2}{n+1}.
\]
\end{theorem}
\begin{proof}
固定 \(0<r<1\)。Taylor 级数在圆周 \(|z|=r\) 上一致收敛，因此有限部分和的平方模积分趋于 \(f\) 的平方模积分。先对有限和使用角频率正交性，再取极限，得到
\[
 \int_0^{2\pi}|f(re^{\mathrm{i}\theta})|^2\,\mathrm{d}m(\theta)
 =\sum_{n=0}^{\infty}|a_n|^2r^{2n}.
\]


对于 \(0<R<1\)，两边乘以 \(2r\)，从 \(0\) 积到 \(R\)。右端各项非负，Tonelli 定理允许交换求和与径向积分，给出
\[
 \int_{|z|<R}|f(z)|^2\,\mathrm{d}A(z)
 =\sum_{n=0}^{\infty}\frac{|a_n|^2}{n+1}R^{2n+2}.
\]
令 \(R\uparrow1\)，左边由区域递增的单调收敛趋于全盘积分；右边也由非负项的单调收敛趋于所述级数。这证明了\eqref{b3:eq:28-bergman-coefficient-norm}，包括无穷值情形。将同一公式用于全纯函数 \(f-p_N\)，便得到精确的尾范数；在级数有限时，这个尾和趋于零。
\end{proof}

反方向也需要说明：任意满足 \(\sum_{n\ge0}|a_n|^2/(n+1)<\infty\) 的复数列，都确实给出圆盘上的全纯函数。对 \(0\le r<1\)，Cauchy--Schwarz 不等式给出
\begin{equation}\label{b3:eq:28-weighted-series-control}
 \begin{aligned}
 \sum_{n=0}^{\infty}|a_n|r^n
 &\le\left(\sum_{n=0}^{\infty}\frac{|a_n|^2}{n+1}\right)^{1/2}
 \left(\sum_{n=0}^{\infty}(n+1)r^{2n}\right)^{1/2}
 \\
 &=\frac1{1-r^2}
 \left(\sum_{n=0}^{\infty}\frac{|a_n|^2}{n+1}\right)^{1/2}.
 \end{aligned}
\end{equation}
最后的求和由几何级数逐项求导得到。对尾级数使用同一估计，可知幂级数在每个闭子圆盘上一致绝对收敛，因而定义全纯函数；再用\cref{b3:thm:28-bergman-coefficients}即可确认它属于 \(A^2\)。

\begin{corollary}[Bergman 标准正交基]\label{b3:cor:28-bergman-basis}
在单位圆盘的归一化面积 Bergman 空间 \(A^2(\mathbb D)\) 中，函数
\[
 e_n(z)=\sqrt{n+1}\,z^n\quad(n=0,1,2,\ldots)
\]
构成标准正交基。若 \(f(z)=\sum_{n\ge0}a_nz^n\in A^2(\mathbb D)\)，则
\[
 \langle f,e_n\rangle_{A^2}=\frac{a_n}{\sqrt{n+1}},
 \quad
 f=\sum_{n=0}^{\infty}\langle f,e_n\rangle_{A^2}e_n,
\]
后一个级数同时在面积范数及局部一致意义下收敛。映射
\[
 f\longmapsto\left(\frac{a_n}{\sqrt{n+1}}\right)_{n\ge0}
\]
是从 \(A^2(\mathbb D)\) 到复序列空间 \(\ell^2(\mathbb N_0)\) 的满等距线性映射。
\end{corollary}
\begin{proof}
Taylor 部分和的面积范数收敛表明，多项式的线性空间在 \(A^2\) 中稠密，故先前的标准正交系是基。对部分和与 \(e_n\) 作内积，有限正交计算给出 \(a_n/\sqrt{n+1}\)；内积连续性允许取部分和的范数极限，得到同一公式。范数恒等式给出等距性；对于任意 \((c_n)\in\ell^2\)，取 \(a_n=\sqrt{n+1}\,c_n\)，前述加权级数构造给出相应的 \(A^2\) 函数，因此映射是满射。局部一致收敛也可由面积范数收敛和\cref{b3:prop:28-area-local-control}直接推出。
\end{proof}

现在每个全纯平方可积函数都由一组平方可和的正交坐标表示。与 Hardy 空间的差别集中在权重 \(1/(n+1)\)：它既说明面积可积容许较大的 Taylor 系数，也决定下一节的核中为何出现二次分母。
